% Champ de gravitation de la Terre : lignes radiales, champ localement uniforme (L'essentiel, chapitre 11)
\begin{tikzpicture}[x=1cm,y=1cm,
  ligne/.style={otGreenDark!80,line width=0.6pt,-{Stealth[length=1.7mm]}}]
\shade[ball color=otNavy!35] (0,0) circle (1.0);
\node[pcLegendeG] at (0,0) {Terre};
\foreach \a in {0,30,...,330}{\draw[ligne] (\a:2.1) -- (\a:1.1);}
\draw[otMuted,line width=0.5pt,dash pattern=on 2pt off 1.5pt] (60:1.0) -- (3.9,1.6);
\draw[otMuted,line width=0.5pt,dash pattern=on 2pt off 1.5pt] (40:1.0) -- (3.9,-0.2);
\begin{scope}[shift={(3.9,-0.2)}]
  \draw[otMuted,line width=0.5pt] (0,0) rectangle (2.6,1.8);
  \fill[otGreen!25] (0,0) rectangle (2.6,0.25);
  \foreach \xx in {0.35,0.85,...,2.4}{\draw[ligne] (\xx,1.65) -- (\xx,0.35);}
  \node[pcLegende,anchor=west,fill=white,inner sep=1pt] at (2.65,1.0) {$\vec{g}$};
  \node[pcLegende,anchor=north,align=center] at (1.3,-0.1) {près du sol : champ\\localement uniforme};
\end{scope}
\end{tikzpicture}
\par\smallskip\legende{Le champ de gravitation terrestre est dirigé vers le centre de la Terre ; à l'échelle d'une pièce, il est uniforme.}
